\chapter{基础数据结构}

\begin{overview}
\begin{itemize}
    \item \texttt{vector}、栈、队列、链表、堆支持的操作，及其实现
    \item 常见的可以使用这些数据结构解决的问题
\end{itemize}
\end{overview}




提到复杂度时默认指最坏复杂度

\section{数组}

相信大家都会使用常数长度的数组。但是也有开变量长度的数组的方法：\verb|int* a=new int[n]|。\graffiti{很多编译器也支持直接定义 \texttt{a[n]}，但这不是标准的。}通过这种方法，我们可以开每维长度不同的二维数组：
\begin{lstlisting}[language=C++]
int pool[1000000]; // 数组充当内存池
int* a[1000]; // 定义一个指针数组充当二维数组
int* cur=pool; // 内存池指针
for (int i=0; i<1000; i++) {
    // a[i] 占用 [cur, cur+L[i]) 这一段空间
    a[i]=cur; 
    cur+=L[i];
}
a[i][j]; // 访问 a[i][j]
\end{lstlisting}
相比 \texttt{vector} 套 \texttt{vector}，这种方法的内存是连续的，且常数更小。

\section{\texttt{vector}}

变长数组。基本可以完全替代数组。

\subsection{基本操作}

\begin{itemize}
    \item 均摊 $O(1)$ 的末尾插入 \texttt{push\_back}
    \item 均摊 $O(1)$ 的末尾删除 \texttt{pop\_back}
    \item $O(1)$ 的随机访问\footnotemark \texttt{a[i]}
    \item $O(\mathrm{size}-i)$ 加均摊 $O(1)$ 在第 $i$ 个位置插入删除
\end{itemize}

参考：https://cppreference.cn/w/cpp/container/vector

\footnotetext{
指任意访问某个 \texttt{a[i]}
}


\subsection{常见问题}
在使用时，需要注意一些陷阱：
\begin{itemize}
    \item 大量 \texttt{push\_back} 会带来很大的常数。可以使用 \texttt{resize} 一次性分配需要的长度，这样和数组基本一样快
    \item 迭代器失效：修改导致长度 \texttt{size} 变化后，迭代器可能变得不可用
    \item \texttt{pop\_back} 不会释放内存，所以占用空间正比于历史最大 \texttt{size} 而不是当前 \texttt{size}
    \item \texttt{clear} 不会释放内存。可以用 \texttt{vector<T>().swap(a)} 来释放 
\end{itemize}

\subsection{实现}
我们给出一种 \texttt{vector} 的实现：如果 \texttt{size} 超过了当前容量，那么把当前容量加倍，重新申请空间，把所有东西都搬过去。

每次重新分配后，容量约等于两倍 \texttt{size}。触发下一次重新分配需要 $\Omega(\mathtt{size})$ 次修改，重新分配的代价是 $O(\mathtt{size})$。所以，平均每次修改的代价是 $O(1)$。

    \begin{question}
        如果额外要求空间线性于 \texttt{size} 而不是历史最大 \texttt{size} 呢？
    \end{question}
    \begin{solution}
    \texttt{size} 小于当前容量的 $1/4$ 时，把容量减半并重新分配空间（为什么是 $1/4$ 而不是 $1/2$？）
    \end{solution}
   
\section{链表}   

\subsection{基本操作}
    链表支持

    \begin{itemize}
        \item $O(1)$ 在指定元素处插入删除单个元素/另一个链表（例如可以实现 $O(1)$ 拼接）
        \item $O(\text{位置})$ 随机访问
        \item $O(\mathtt{size})$ 的空间
    \end{itemize}

    \begin{example}
        洛谷P10466 邻值查找 / 洛谷P1168 中位数
    \end{example}

    插入往往需要平衡树之类的数据结构，但是如果倒着来就只需要先排序然后一直删除。如果排序可以线性，那么就能做到线性。

\begin{remark}
    值域为 $n^{O(1)}$ 的排序可以通过\textbf{多关键字计数排序}，或者说\textbf{基数排序}来做到线性。
\end{remark}

\section{栈和队列}
    基本操作是广为人知的，不再赘述。显然，栈在功能上弱于 \texttt{vector}，因此可以用 \texttt{vector} 实现栈。而队列则不能直接用 \texttt{vector} 实现。如果队列大小已知不超过某个上界，那么可以用数组模拟循环队列。而如果想做到 $O(\mathtt{size})$ 空间，就需要用链表或者双端队列实现

\section{双端队列}
    双端队列在功能上强于 \texttt{vector} 和队列，但是实现更复杂，常数也更大。双端队列的支持的操作如下：
    \begin{itemize}
        \item 均摊 $O(1)$ 头尾插入删除
        \item $O(1)$ 随机访问
    \end{itemize}
    注意如果没有随机访问的要求，那么可以直接用链表完成。

    STL 的 \texttt{deque} 即是标准的双端队列。然而，需要开很多双端队列时，要谨慎使用 STL 的 \texttt{deque}，因为一个空 \texttt{deque} 可能占用几百到几千字节的空间。\graffiti{NOI 2022 时有许多人因为 MLE 失去了一整道题的分数。} 

    更好的实现是对顶栈。设想有两个栈，其底部粘在一起。于是其整体就成为一个两端开口的双端队列。唯一的问题出在某个栈空了还要继续删的情况。此时，我们暴力把整个序列倒出来重构即可。换句话说，我们重新把序列均分到两个栈里。

    使用 \texttt{vector} 来实现栈就能支持随机访问了。如果需要很紧的空间，可以用之前提到的 $O(\mathtt{size})$ 空间的 \texttt{vector}。
    
    每次重构后，两个 \texttt{vector} 都有大概 $\mathtt{size}/2$ 个元素。到下次重构之前，至少有 $\mathtt{size}/2+|\mathtt{size}'-\mathtt{size}/2|=\Omega(\mathtt{size}')$ 次操作。重构花费 $O(\mathtt{size}')$ 的代价，所以插入删除是均摊 $O(1)$ 的。

\subsection{维护信息}
    在两个 \texttt{vector} 中各自维护自己的前缀信息，这样可以 $O(1)$ 查询当前双端队列中跨过分断点的任意区间的分治信息。\graffiti{分治信息的定义和例子将在下一章中给出。}
    \begin{example}
        支持双端队列操作、每次询问以体积 $V_i$ 做 01 背包的结果。需要 $O((n+q)V)$。
    \end{example}

    假设当前询问区间是 $[l,r]$，那么通过双端队列操作，我们可以 $O(1)$ 移动到 $[l\pm 1,r\pm 1]$。这样，我们就可以通过移动区间端点来在 $\sum_i |l_i-l_{i-1}|+|r_i-r_{i-1}|$ 的时间内回答多个区间询问。在很多问题中，这个总代价往往可以是线性的。

    \begin{example}[滑动窗口]
        给定序列 $a$ 和整数 $k$，查询所有 $[i,i+k-1]$ 的分治信息（如最大子段和）。
    \end{example}

    \begin{example}[“不删除双指针”]
        询问最短的 $\gcd>1$ 的区间。$\gcd>1$ 可以换成 $\max\ge x$、按位或 $\ge x$ 之类的有单调性的条件。
    \end{example}

\section{二叉堆}

\subsection{基本操作}

    \begin{itemize}
        \item $O(1)$ 查询最小值
        \item $O(\log n)$ 插入
        \item $O(\log n)$ 删除
        \item $O(\log n)$ Decrease-Key
        \item $O(n)$ 合并
    \end{itemize}

\begin{remark}
    理论上最好的堆——Brodal Queue 支持
    \begin{itemize}
        \item $O(1)$ 查询最小值
        \item $O(1)$ 插入
        \item $O(\log n)$ 删除
        \item $O(1)$ Decrease-Key
        \item $O(1)$ 合并
    \end{itemize}

    但是实现复杂、常数很大。

    算法/数据结构总是在复杂度、实际表现、实现难度、常数、功能之间做取舍。

    还有其他一些常见的堆：左偏树、斜堆、随机堆、配对堆、二项堆、斐波那契堆。
\end{remark}

\section{扩展类问题}
    \graffiti{可惜没时间讲。}
    从若干个初始状态开始，每个状态可以扩展出若干个后继状态，保证扩展出来的状态的权值小于原状态的权值。每次取权值最小的状态扩展，重复若干次。

    通过一定的设计，很多求前 $k$ 大/小的问题可以变成这类问题。使用堆来维护所有目前的状态，每次取出最小的一个来扩展，注意去重。扩展形成一张图，复杂度为 $O(M\log M)$ 或 $O(k\log k+M)$（取决于堆），其中 $M$ 为前 $k$ 小的点邻接的\textbf{边数}。

    如果每个点的扩展方式都是相同的，并且只有常数种扩展方式，并且扩展出来的状态的权值关于原状态的权值是单调的，那么可以使用队列替换堆来做到 $O(M)$。
\begin{example}[序列合并]

\end{example}
\begin{solution}
        方法一：$(i,j)$ 扩展出 $(i,j+1)$，初始状态为所有 $(i,1)$

    方法二：$(i,j)$ 扩展出 $(i+1,j)$ 和 $(i,j+1)$，但是要去重

    方法三：https://www.cnblogs.com/zkyJuruo/p/18428124
\end{solution}

\begin{example}[蚯蚓]
    
\end{example}
\begin{solution}
    首先处理掉“除了新的以外都加 $q$”，这可以变成全局加 $q$ 然后新的减掉 $q$。

    这样就是 $x$ 扩展出 $\lfloor p(x+tq)\rfloor-(t+1)q$ 和 $(x+tq)-\lfloor p(x+tq)\rfloor-(t+1)q$

    符合我们用队列替换堆时的要求
\end{solution}
\begin{example}[常数种非负边权的最短路]
    BFS 扩展时，从 $x$ 扩展出 $O(1)$ 种固定的新状态

    可以使用队列替换堆。
\end{example}
\begin{example}[前 $k$ 小子集和]

\end{example}
\begin{solution}
    考虑一个 DFS 过程，枚举下一个选择的数在哪

    然而这样会扩展出 $\Theta(n)$ 个状态

    多叉树转二叉树：左儿子右兄弟（注意兄弟的顺序）
\end{solution}
\begin{example}[限制大小的前 $k$ 小子集和]
    要求集合大小在 $[L,R]$ 内。
\end{example}
\begin{solution}
    考虑怎么枚举 $\mathtt{size}=m$ 的集合。假设选了 $a_{i_1},\ldots,a_{i_m}$，转为枚举 $d_j=i_j-i_{j-1}$，只需要 $\sum_j d_j\le n-1$。
    
    初始所有 $d_j=1$，每次枚举下一个需要增加的 $d_j$ 即可。

    同样用儿子兄弟表示法转为二叉树
\end{solution}
\begin{remark}
    这类问题应该还有很大研究空间，现在知道的有 \href{https://www.cnblogs.com/AK-IOI/p/18418766}{这篇} 和 \href{https://www.cnblogs.com/Cry-For-theMoon/p/17629760.html}{这篇}。有兴趣的同学可以自行研究或者和我一起讨论一下
\end{remark}
\begin{example}[合并果子]
    由于合并的结果总是越来越大的，可以使用队列替换堆。
\end{example}
\begin{remark}
    这就是 Huffman 编码。我们有结论

    \begin{itemize}
        \item Huffman 编码可以最小化 $\sum_i \mathrm{dep}_i\cdot w_i$
        \item $w_i$ 的深度大约为 $\log\frac{W}{w_i}$，其中 $W=\sum w_i$
    \end{itemize}

    在\textbf{不均匀分治}（如树剖、点分治、分治FFT等）中，可以使用 Huffman 树来降低复杂度或常数，以后会提到。
\end{remark}

\section{单调队列和单调栈}

令 $b_i=\max\{a_i,a_{i+1},\ldots,a_r\}$，即 $b$ 是只考虑 $a_1,\ldots,a_r$ 时的 $a$ 的后缀 $\max$。我们可以在 $r$ 从 $1$ 扫到 $n$ 的过程中均摊 $O(1)$ 地隐式维护所有 $b_i$。
具体来说，维护一个栈，栈中元素是 $a$ 的一个子序列。假设自底向上的第 $k$ 个元素是 $a_{i_k}$，那么对任何 $i_{k-1}<j\le i_k$，都有 $b_j=a_{i_k}$，也就是说栈中每个元素都表示一段以它为结尾的区间的 $b_i$。由于 $b$ 是不升的，栈中元素也是不升的。当 $r$ 移动到 $r+1$ 时，考虑 $a_{r+1}$ 的影响。假设当前栈中有 $k$ 个元素。如果 $a_{r+1}<a_{i_k}$，那么 $a_{r+1}$ 不会对 $b_{[1,r]}$ 造成任何影响，只需要把 $a_{r+1}$ 加入栈中表示 $b_{r+1}$ 即可。如果 $a_{r+1}\ge a_{i_k}$，那么我们可以想象 $b_{[1,r]}$ 发生的变化：一个后缀会全都变成 $a_{r+1}$。在栈中，这表现为弹出若干栈顶，直到栈顶的 $a_{i_k}>a_{r+1}$，最后再把 $a_{r+1}$ 加入栈中。暴力做这个过程的均摊复杂度就是正确的，因为每个 $a_i$ 最多入栈出栈各一次。

\begin{exercise}
    把 $\max$ 换成 $\min$。
\end{exercise}

    \begin{example}[询问端点单调 RMQ]\label{example:sliding-rmq}
    多次查询 $\max\{a_{l_i},\ldots,a_{r_i}\}$，其中 $l_i,r_i$ 都单调不降。可以强制在线（指 $a$ 可能依赖先前的询问结果）

    虽然双端队列也可以做，但是这种问题单调队列更方便。
    \end{example}
    \begin{solution}
        单调栈中维护了所有的后缀 $\max$，答案就是下标 $\ge l_i$ 的第一个元素。由于 $l_i$ 不降，我们直接弹出队头就可以均摊 $O(1)$ 地维护这个事情。
    \end{solution}
    \begin{exercise}[滑动窗口]
        对每个 $1\le l\le n-k+1$ 查询 $\max\{a_i\mid i\in [l,l+k-1]\}$。
    \end{exercise}
    \begin{exercise}\label{exercise:monotone-stack-rmq}
        在\cref{example:sliding-rmq}的基础上减弱问题的假设。
        \begin{enumerate}
            \item 去掉 $l_i$ 单调的假设，允许询问复杂度变大一些
            \item 再通过离线去除 $r_i$ 单调的假设
            \item *再通过可持久化数组支持在线
            \item 在 $n\le 64$ 时，通过位运算来存储每个 $r$ 的栈中元素集合，从而实现 $O(n)-O(1)$ 的在线 RMQ。我们将在下一章中配合其他技巧做到任意 $n$ 的线性 RMQ。
        \end{enumerate}
    \end{exercise}

    \begin{example}[求前驱后继]    
    对每个元素，查询它左/右边第一个比它大/小的数的位置
    \end{example}
    \begin{solution}
        右边就是把序列倒过来的左边，所以只考虑左边。这就是单调栈中倒数第二个元素。
    \end{solution}

    \begin{example}[稍复杂一些的信息维护]
        对每个 $r$，求 $\sum_{l=1}^r (r-l+1)\cdot \max\{a_l,\ldots,a_r\}$。强制在线。（提示：要记得单调栈中每个元素都代表了一段区间的 $\max\{a_l,\ldots,a_r\}$。）
    \end{example}
    \begin{solution}
        假设栈中元素是 $a_{i_1}, a_{i_2}, \ldots, a_{i_k}$，其中 $i_1 < i_2 < \ldots < i_k$，那么答案是
        $$
        \begin{aligned}
        \sum_{j=1}^k a_{i_j}\sum_{i_{j-1}<l\le i_j}(r-l+1)&=\sum_{j=1}^k a_{i_j}\cdot(i_j-i_{j-1})\cdot(r-i_j+1)\\
        &=(r+1)\sum_{j=1}^k a_{i_j}-\sum_{j=1}^k a_{i_j}(S_{i_j}-S_{i_{j-1}}),
        \end{aligned}
        $$
        其中 $S_i=\sum_{k=1}^i k$。直接维护这个式子即可。
    \end{solution}
    \begin{solution}[离线的另解]
        如果允许离线，那么可以考虑换种方式计算答案：对每个 $a_i$，看看它是哪些区间 $[l,r]$ 的 $\max$。这可以由先前提到的前驱后继来做。
    \end{solution}

    \begin{example}[最值版本]
        求 $\max_{l,r}(r-l+1)\cdot\max\{a_l,\ldots,a_r\}$。
    \end{example}
    \begin{solution}
        使用离线算贡献的方法是最容易的

    在线对每个 $r$ 算的话需要维护凸包。
    \end{solution}

\begin{example}[悬线法]
    有些题给了个二维黑白网格，要求全黑的矩形/正方形的最大面积/面积和之类的东西

    对每个格子计算每个方向全黑能延伸到哪，然后就变成类似例三的问题
\end{example}

\begin{exercise}{困难的例子}
    给定两个序列 $a,b$，求有多少个区间 $[l,r]$ 满足 $\max a[l,r]=\max b[l,r]$。
\end{exercise}
\begin{solution}
    $a,b$ 交替形成一个新序列，在这个新序列上扫的时候用单调栈维护所有后缀 $\max$。在维护这个交替序列的后缀 $\max$ 的同时，我们可以顺便维护出每个后缀 $\max$ 是否同时也是只考虑 $a/b$ 的后缀 $\max$，以及只考虑 $a/b$ 时何时达到此后缀 $\max$。之后容易计算单调栈中每个项的贡献：假设只考虑 $a/b$ 时在原序列的 $l_a/l_b$ 位置达到此后缀 $\max$，则这一项的贡献即为 $\min(l_a,l_b)-\mathrm{pos}$，其中 $\mathrm{pos}$ 为此项在原序列中的位置，另外在此项是由 $b$ 产生的时候 $pos$ 需要再加一。

    也可以扩展到多个序列。
\end{solution}


% =========================================================

% =========================================================

\begin{summary}
\begin{itemize}
    \item 熟知讲过的几种线性数据结构的功能，能够把问题与数据结构进行对应
    \item 扩展类问题的思路
\end{itemize}
\end{summary}
